\documentclass[12pt,a4paper]{article}

% Packages
\usepackage[utf8]{inputenc}
\usepackage[T1]{fontenc}
\usepackage{amsmath,amssymb,amsthm}
\usepackage{graphicx}
\usepackage{hyperref}
\usepackage{booktabs}
\usepackage{geometry}
\usepackage{xcolor}
\usepackage{caption}
\usepackage{subcaption}
\usepackage{enumitem}
\usepackage{url}
\usepackage{doi}

\geometry{margin=1in}

% Theorem environments
\newtheorem{definition}{Definition}
\newtheorem{proposition}{Proposition}
\newtheorem{remark}{Remark}

% Title information
\title{\textbf{Temporal Emergence from Dynamical Observer Coupling:\\
       A Matrix Model for Observer-Dependent Symmetry Breaking}\\[0.5cm]
       \large Version 6: Trajectory-Mean Analysis with Corrected Sampler}

\author{Sutipong Chanpengpad\\
\small Independent Researcher, Chiang Rai, Thailand\\
\small \texttt{dhammawatthumpra@gmail.com}\\
\small \href{https://orcid.org/0009-0001-4069-8576}{ORCID: 0009-0001-4069-8576}}

\date{\today}

\begin{document}

\maketitle

\begin{abstract}
The Structural Time Framework (STF) proposes that time emerges from 
observer-dependent ordering of pre-temporal states. We present a 
computational demonstration using a matrix model where an explicit 
observer system $Y$ dynamically creates the temporal direction in 
the main system $X$, without any external parameters.

Our model employs a directional coupling mechanism where the observer 
generates a direction vector $v$ from its matrix traces, and dimensions 
of $X$ aligned with $v$ receive a temporal expansion term 
$-g_{XY} \hat{v}_\mu^2 \text{Tr}(X_\mu^4)$. All results reported here use the
\textbf{v6 trajectory-mean estimator}: observables are averaged over the
measurement window (30 recorded sweeps) rather than read from a single final
configuration. Monte Carlo simulations with statistical analysis (5 seeds per
configuration, with a $30$-seed follow-up for the observer-dimension study)
show:

\begin{enumerate}[leftmargin=*]
\item Temporal emergence occurs robustly for observer dimensions
      $d \in \{2,3,4,5\}$ and system sizes $N \in \{4,5,\ldots,16\}$, with
      mean ratios clustering in $4.4$--$4.8$ and alignment at $100\%$ in every
      configuration.
\item No crossover or two-regime scaling is observed: emergence remains
      healthy across the entire tested range of $N$ up to $N=16$, with no
      threshold at which fixed coupling $g_{XY}=0.8$ fails to temporalize the
      system.
\item Perfect alignment between the temporal direction $X_{\max}$ and 
      observer direction $v_{\max}$ in all healthy emergence cases.
\item A follow-up with $n=30$ seeds per condition shows a weak but
      detectable monotonic decrease of the mean ratio with observer
      dimension (slope $\approx -0.08$ per unit $d$, $p=0.003$, $R^2=0.07$),
      although the pairwise comparison between $d=3$ and $d=4$ is not
      significant at this sample size ($t=0.52$, $p=0.61$). The earlier
      $n=5$ result ($p=0.049$) was a small-sample fluctuation.
\end{enumerate}

\textbf{Important Caveats:} This work represents a toy-level 
mathematical exploration using small matrix sizes ($N \leq 16$). 
Results demonstrate mathematical consistency, not physical reality.

\medskip
\noindent\textbf{Keywords:} dynamical observer, temporal emergence, 
matrix models, symmetry breaking, computational demonstration
\end{abstract}

\tableofcontents
\newpage

%=============================================================================
\section{Introduction}\label{sec:intro}
%=============================================================================

The nature of time has been debated by philosophers and physicists for 
centuries~\cite{page_wootters_1983, barbour_1999, rovelli_2018}. The Structural Time Framework (STF)~\cite{chanpengpad_stf} 
proposes that time is not fundamental but emerges from observer-dependent 
ordering of pre-temporal states.

In our previous work~\cite{chanpengpad_sgoed_v2}, we presented a toy 
model where temporal emergence was driven by an external adapter parameter 
$\kappa_2$. While successful, this approach had a fundamental limitation: 
the observer was represented as a parameter, not as a dynamical system.

This paper addresses that limitation by introducing an \textbf{explicit 
dynamical observer}---a matrix system $Y$ that dynamically creates the 
temporal direction in the main system $X$ through coupling, without any 
external parameters.

\textbf{Important Disclaimer:} We emphasize that this work is a 
\textbf{toy-level computational demonstration}, not empirical physics. 
Our simulations use small matrix sizes ($N \leq 16$), finite Monte Carlo 
samples, and idealized actions. The results demonstrate mathematical 
consistency, not physical reality.

%=============================================================================
\section{Dynamical Observer Model}\label{sec:model}
%=============================================================================

We consider two coupled matrix systems:

\textbf{Main System $X$:} $D$ Hermitian matrices $X^\mu$ of size $N \times N$.

\textbf{Observer System $Y$:} $d$ Hermitian matrices $Y^a$ of size $N \times N$.

The total action is:
\begin{equation}
S_{\text{total}} = S_X + S_Y + S_{\text{coupling}}
\end{equation}

where $S_X$ and $S_Y$ are IKKT-like actions with stability gates, and the 
directional coupling is:
\begin{equation}
S_{\text{coupling}} = -g_{XY} \sum_\mu \hat{v}_\mu^2 \text{Tr}(X_\mu^4)
\end{equation}

The observer generates direction $\hat{v}$ from its traces. This coupling 
temporalizes dimensions aligned with $\hat{v}$.

\subsection{Theoretical Motivation for the Coupling}

The coupling term can be understood from three perspectives:

\begin{enumerate}[leftmargin=*]
\item \textbf{Effective Field Theory:} Leading-order term in $1/N$ expansion.
\item \textbf{Symmetry Breaking:} Explicitly breaks SO(D) $\rightarrow$ SO(D-1).
\item \textbf{Measurement Analogy:} Observer selects direction, system responds.
\end{enumerate}

\subsection{Why the Observer Does Not Temporalize Itself}

The distinction between observer and system is relational, not fundamental. 
The asymmetry in coupling reflects the measurement paradigm: the apparatus 
measures the system, not vice versa. This avoids the self-observation paradox 
analogous to Wigner's friend.

%=============================================================================
\section{Computational Methodology}\label{sec:method}
%=============================================================================

We use Metropolis-Hastings sampling~\cite{metropolis_1953} with:
\begin{itemize}
\item Matrix size: $N \in \{4, 5, 6, 7, 8\}$, with a larger-$N$ extension to
      $N \in \{9, 10, 12, 14, 16\}$ (Section~\ref{sec:fss})
\item System dimensions: $D = 6$
\item Observer dimensions: $d \in \{2, 3, 4, 5\}$
\item Coupling strength: $g_{XY} \in [0.8, 1.15]$
\item Thermalization: 20 sweeps
\item Measurement: 30 sweeps
\item Multiple seeds: $\{42, 43, 44, 45, 46\}$ (base set), plus
      $\{100,\dots,129\}$ for the $30$-seed observer-dimension study
\end{itemize}

\textbf{Thermalization check:} the observables reported below are robust to
increasing the thermalization budget. Re-running representative configurations
with $n_{\mathrm{therm}} = 20, 40, 80, 160$ leaves the mean ratio and
alignment unchanged (e.g.\ $R \approx 4.72$--$4.84$ at $N=8$, $d=3$,
$g_{XY}=0.8$), and the wall fraction stays at zero throughout. We therefore
regard $n_{\mathrm{therm}} = 20$ as sufficient for the healthy (v6) regime
studied here. (For the separate two-way-coupling model of v7, in which the
observer itself can be driven into the stability gate, $n_{\mathrm{therm}}=20$
was found insufficient and a larger budget is required; that model is not
reported in this paper.)

The sampler is the \textbf{correctness-first v6 engine}
(\texttt{code/sgoed\_core\_v6.py}). It computes Metropolis acceptance via an
incremental local-energy delta rather than a full-action recomputation; the
delta was validated against the full action to machine precision (maximum
relative error $< 10^{-9}$ over 4000 randomized trials, including wall-regime
edge cases). The v6 sampler updates every matrix element at every sweep. This
replaces the earlier v5 engine, which applied a \texttt{step=2} element-skip
optimization for $N>6$ (updating only even-indexed matrix elements)---a
shortcut that we subsequently identified as the source of a spurious
finite-size crossover and a spurious variance spike at $N=8$ (Section~\ref{sec:absence_crossover}).

All observables are recorded over the measurement window and then averaged
(\textbf{trajectory mean}), rather than read from a single final
configuration.

Observables measured, with $\text{extent}(X_\mu) \equiv \text{Tr}(X_\mu^2)/N$:
\begin{align}
\text{Ratio} &= \frac{\text{extent}(X_{\max})}{\text{mean}(\text{extent}(X_{\neq \max}))} \\
\text{Alignment} &= \mathbb{I}(\arg\max_\mu \text{extent}(X_\mu) == \arg\max_\mu \hat{v}_\mu^2) \\
H &= -\sum_\mu p_\mu \ln p_\mu, \quad p_\mu = \frac{\text{extent}(X_\mu)}{\sum_\nu \text{extent}(X_\nu)}
\end{align}

A configuration is classified as \textbf{Healthy} if the mean Ratio over
the sampled seeds satisfies $1.5 < \overline{R} < 10.0$; configurations with
$\overline{R} \leq 1.5$ are labeled Weak. The upper bound $R=10$ is chosen to
match the stability-gate cutoff on the per-matrix \emph{extent}
(\texttt{max\_extent}$=10.0$) used in the action itself. We emphasize that
the two quantities are distinct: the gate acts on $\text{extent}(X_\mu) =
\text{Tr}(X_\mu^2)/N$ of a \emph{single} matrix (a dimensionful scale), whereas
$R$ is the \emph{dimensionless ratio} of the largest extent to the mean of the
remaining extents. When the aligned dimension is driven toward the gate
ceiling, its extent saturates just below $10.0$ while the others remain near
$1$, so $R$ approaches $\sim 10$ even though no individual extent ever exceeds
the cutoff (the gate penalty suppresses growth before the threshold is
crossed). None of the configurations reported below have an individual extent
exceeding $10.0$.

%=============================================================================
\section{Results}\label{sec:results}
%=============================================================================

\subsection{Statistical Robustness}

Table~\ref{tab:stats} summarizes results for optimal parameters.

\begin{table}[ht]
\centering
\caption{Statistical analysis ($N=6$, $g_{XY}=0.8$)}
\label{tab:stats}
\begin{tabular}{ccccc}
\toprule
$d$ & Mean Ratio & Std Dev & CV (\%) & Alignment \\
\midrule
3 & 4.65 & $\pm$0.23 & 4.9\% & 100\% (30/30) \\
4 & 4.61 & $\pm$0.32 & 7.0\% & 100\% (30/30) \\
\bottomrule
\end{tabular}
\end{table}

A follow-up with $n=30$ seeds per condition (rather than the initial $n=5$)
gives a paired t-test of $t = 0.52$, $p = 0.61$ for the $d=3$ versus $d=4$
comparison. This is not significant. The marginally significant result
reported in an earlier draft of this version ($t=2.804$, $p=0.049$ at
$n=5$) is therefore attributed to a small-sample fluctuation and is
withdrawn: with a larger sample the pairwise difference between $d=3$ and
$d=4$ is not resolved.

Across all four observer dimensions ($d=2,3,4,5$, $n=30$ each), however, the
mean ratio decreases monotonically from $4.74$ ($d=2$) to $4.49$ ($d=5$), a
linear trend of slope $\approx -0.08$ per unit $d$ ($p=0.003$,
$R^2=0.07$). The trend is weak---it explains only $\sim 7\%$ of the
seed-to-seed variance---but is statistically detectable at this sample size,
consistent with a mild ``dimensional dilution'' of the directional coupling.
It should not be read as evidence that any single pair of dimensions
differs.

%=============================================================================
\subsection{Observer-Dimension Phase Diagram}\label{sec:phase_diagram}
%=============================================================================

To investigate the role of observer complexity, we varied the observer
dimension $d \in \{2,3,4,5\}$ while keeping $N=6$, $D=6$, and
$g_{XY}=0.8$ fixed. Each configuration was evaluated over five independent
random seeds.

\begin{table}[ht]
\centering
\caption{Phase diagram by observer dimension ($N=6$, $g_{XY}=0.8$)}
\label{tab:phase_diagram}
\begin{tabular}{cccccc}
\toprule
$d$ & Mean Ratio & Std Dev & CV (\%) & Alignment & Status \\
\midrule
2 & 4.74 & $\pm$0.15 & 3.1\% & 100\% (30/30) & Healthy \\
3 & 4.65 & $\pm$0.23 & 4.9\% & 100\% (30/30) & Healthy \\
4 & 4.61 & $\pm$0.32 & 7.0\% & 100\% (30/30) & Healthy \\
5 & 4.49 & $\pm$0.46 & 10.3\% & 100\% (30/30) & Healthy \\
\bottomrule
\end{tabular}
\end{table}

Several features are noteworthy.

\begin{enumerate}[leftmargin=*]
\item \textbf{No strict threshold at $d=3$:}
Contrary to our initial expectation, even $d=2$ produces healthy temporal
emergence with perfect alignment. This suggests that the minimum observer
complexity required for temporalization is lower than anticipated.

\item \textbf{Robust emergence, with a weak downward trend in $d$:}
Mean ratios cluster in a narrow range ($4.49$--$4.74$) across all four tested
observer dimensions, and alignment is perfect (100\%, 30/30 seeds) in every
case. With $n=30$ seeds per dimension the mean ratio decreases monotonically
from $4.74$ ($d=2$) to $4.49$ ($d=5$), a weak trend (slope $\approx -0.08$
per unit $d$, $p=0.003$, $R^2=0.07$; Table~\ref{tab:stats}); the paired
comparison between $d=3$ and $d=4$ is, however, not significant at this
sample size ($p=0.61$). We therefore report a mild ``dimensional dilution''
of the directional coupling rather than a robust pairwise difference, and we
find no evidence for an optimal or ``sweet spot'' observer dimension in this
range. (An earlier v5 draft of this table reported a spurious $d=4$ ``sweet
spot'' and reduced $d=5$ alignment; those numbers could not be reproduced
from the corrected v6 sampler and have been replaced with freshly re-run,
seed-logged values.)
\end{enumerate}

%=============================================================================
\subsection{Finite-Size Scaling and Parameter Sensitivity}\label{sec:fss}
%=============================================================================

We next tested whether the observed temporal emergence is a finite-size
artifact by varying the matrix size $N$. We first fixed $d=3$ and
$g_{XY}=0.8$ and examined $N \in \{4,5,6,7,8\}$.

\begin{table}[ht]
\centering
\caption{Finite-size scaling with fixed $g_{XY}=0.8$ ($d=3$)}
\label{tab:fss_fixed}
\begin{tabular}{cccccc}
\toprule
$N$ & Mean Ratio & Std Dev & CV (\%) & Alignment & Status \\
\midrule
4 & 4.63 & $\pm$0.19 & 4.2\% & 100\% (5/5) & Healthy \\
5 & 4.62 & $\pm$0.18 & 4.0\%  & 100\% (5/5) & Healthy \\
6 & 4.75 & $\pm$0.17 & 3.5\%  & 100\% (5/5) & Healthy \\
7 & 4.75 & $\pm$0.12 & 2.6\%  & 100\% (5/5) & Healthy \\
8 & 4.72 & $\pm$0.05 & 1.0\%  & 100\% (5/5) & Healthy \\
\bottomrule
\end{tabular}
\end{table}

To test whether this robustness is a small-$N$ accident, we extended the scan
to $N \in \{9, 10, 12, 14, 16\}$ (same $d=3$, $g_{XY}=0.8$, five seeds each):

\begin{table}[ht]
\centering
\caption{Larger-$N$ extension with fixed $g_{XY}=0.8$ ($d=3$)}
\label{tab:fss_large_n}
\begin{tabular}{cccccc}
\toprule
$N$ & Mean Ratio & Std Dev & CV (\%) & Alignment & $\tau_{\mathrm{int}}$ \\
\midrule
9  & 4.75 & $\pm$0.10 & 2.0\% & 100\% (5/5) & 1.8 \\
10 & 4.77 & $\pm$0.06 & 1.3\% & 100\% (5/5) & 2.2 \\
12 & 4.76 & $\pm$0.10 & 2.1\% & 100\% (5/5) & 1.8 \\
14 & 4.75 & $\pm$0.03 & 0.6\% & 100\% (5/5) & 1.6 \\
16 & 4.72 & $\pm$0.07 & 1.6\% & 100\% (5/5) & 1.9 \\
\bottomrule
\end{tabular}
\end{table}

The robustness extends to $N=16$: the ratio remains at $R \approx 4.7$--$4.8$,
alignment stays perfect, and the wall fraction is zero throughout. The
integrated autocorrelation time $\tau_{\mathrm{int}} \approx 1.6$--$2.2$
(measurement sweeps) confirms that the 30 recorded trajectory points are only
weakly correlated, so the trajectory-mean estimates are well behaved.

In contrast to an earlier v5 analysis, the corrected sampler shows
\textbf{no finite-size crossover}: the ratio remains healthy and statistically
stable across the entire tested range $N=4$--$16$.

\begin{enumerate}[leftmargin=*]
\item \textbf{No crossover regime:}
At every tested size the mean ratio stays within $4.6$--$4.8$ and alignment
is perfect (100\%). There is no $N$ at which fixed coupling
$g_{XY}=0.8$ fails to produce temporalization.

\item \textbf{CV decreases monotonically with $N$:}
The coefficient of variation falls steadily from $4.2\%$ at $N=4$ to about
$1\%$ at $N=8$ and remains at or below $\sim 2\%$ through $N=16$. Larger
systems produce more tightly self-averaging measurements, consistent with the
trajectory-mean estimator averaging over more matrix degrees of freedom.
\end{enumerate}

The apparent crossover at $N\approx 7$ reported in v5 (ratio dropping to
$\approx 1.8$ with reduced alignment) is not reproduced here. As documented
in Section~\ref{sec:absence_crossover}, it was an artifact of the v5
\texttt{step=2} element-skip shortcut, not a feature of the model.

%=============================================================================
\subsection{Parameter Tuning at $N=8$}\label{sec:n8_tuning}
%=============================================================================

To examine how the coupling strength affects the strong-coupling side of the
model, we repeated the $N=8$ simulations while varying $g_{XY}$.

\begin{table}[ht]
\centering
\caption{Parameter tuning at $N=8$ ($d=3$)}
\label{tab:n8_tuning}
\begin{tabular}{cccccc}
\toprule
$g_{XY}$ & Mean Ratio & Std Dev & CV (\%) & Alignment & Status \\
\midrule
0.80 & 4.72 & $\pm$0.05 & 1.0\% & 100\% (5/5) & Healthy \\
1.05 & 10.09 & $\pm$0.20 & 2.0\% & 100\% (5/5) & Gate ceiling \\
1.10 & 10.22 & $\pm$0.04 & 0.4\% & 100\% (5/5) & Gate ceiling \\
1.15 & 10.34 & $\pm$0.23 & 2.3\% & 100\% (5/5) & Gate ceiling \\
\bottomrule
\end{tabular}
\end{table}

The corrected sampler changes the picture materially. At $g_{XY}=0.80$ the
system is already healthy ($R \approx 4.72$, CV $1.0\%$), so \textbf{no
parameter tuning is required at $N=8$}. For $g_{XY} \geq 1.05$ the ratio
saturates near $R \approx 10$ with CV near $1$--$2\%$; this is the gate
ceiling described in Section~\ref{sec:method}: the aligned dimension's
\emph{extent} is driven to just below the \texttt{max\_extent}$=10.0$ cutoff
(individual extents peak at $\approx 9.99$ and never exceed $10.0$), so the
dimensionless ratio approaches $\sim 10$ while the gate penalty suppresses
further growth.

Two points follow. First, the non-monotonic CV spike at $g_{XY}=1.05$
(CV $=59.0\%$) reported in v5 is not reproduced here; CV is flat and small
($0.4$--$2.3\%$) across the scanned range. Second, $g_{XY}=0.80$ is the
preferred working value at $N=8$: it produces healthy emergence well inside
the Healthy window, and increasing $g_{XY}$ beyond $\sim 1.05$ only pins the
aligned extent against the gate without further benefit.

%=============================================================================
\subsection{On the Absence of a Crossover Regime}\label{sec:absence_crossover}
%=============================================================================

Unlike our earlier v5 results (which used a final-snapshot estimator and a
\texttt{step=2} element-skip shortcut for $N>6$), the v6 trajectory-mean
results show \textbf{no evidence of a crossover or phase transition} in the
range $N=4$--$16$. The ratio remains stable at $R \approx 4.6$--$4.8$ across
all tested system sizes, with alignment $100\%$ and CV \emph{decreasing} with
$N$. This demonstrates that the apparent crossover at $N \approx 7$---and the
accompanying ``two-regime scaling'' with a tuned optimal coupling
$g_{XY}^{\mathrm{opt}}(N)$---was an artifact of the v5 sampling strategy, not
a feature of the model.

Specifically, the v5 engine applied
\begin{verbatim}
    step = 2 if N > 6 else 1
\end{verbatim}
which caused the X-matrices at $N=7,8$ to have only their even-indexed
elements updated, leaving the odd-indexed elements frozen near their
initial values and preventing the system from thermalizing. The Y-matrices
were not subject to this skip, so the two subsystems were treated
asymmetrically. The v6 engine removes this shortcut and updates every element
at every sweep; once the system is allowed to thermalize fully, the
finite-size ``crossover'' disappears and the mechanism is seen to be robust
across the entire tested range.

We therefore withdraw the earlier two-regime scaling picture and the
power-law estimate $\alpha \approx 0.46$--$1.10$ for
$g_{XY}^{\mathrm{opt}}(N)$. There is no longer a set of tuned optima from
which such an exponent could be computed: $g_{XY}=0.8$ suffices for every
$N \in \{4,5,\ldots,16\}$.

%=============================================================================
\subsection{Physical Interpretation: Robustness and Gate Saturation}\label{sec:crossover_interpretation}
%=============================================================================

With the crossover removed, the two physically meaningful observations are
the \emph{robustness} of temporalization across $N$ and $d$, and the
\emph{saturation} of the aligned extent against the stability gate at strong
coupling.

\begin{enumerate}[leftmargin=*]
\item \textbf{Robustness across system size:}
The mechanism produces a healthy ratio and perfect alignment at every tested
size, with no required change in coupling. This is consistent with the
coupling energy remaining sufficient relative to thermal fluctuations
throughout $N \leq 8$; the decrease of CV with $N$ further indicates that
larger systems self-average, rather than destabilizing.

\item \textbf{Gate saturation at strong coupling:}
For $g_{XY} \gtrsim 1.05$ at $N=8$, the aligned dimension's extent is driven
to the stability-gate ceiling ($\approx 9.99$, just below
\texttt{max\_extent}$=10.0$). The gate penalty term
$10.0\,(\text{extent} - 10.0)^2$ then balances the coupling, pinning the
extent just below the cutoff. The resulting ratio approaches $\sim 10$ while
no individual extent ever exceeds $10.0$; this is the gate doing its job,
not a new physical phase.

\item \textbf{Spectral mechanism (rank-1 condensation):}
The directional coupling acts on the fourth moment
$\mathrm{Tr}(X_\mu^4)$, which strongly favors states in which a single
eigen-direction dominates. At $g_{XY}=0.8,\ N=8$, the selected matrix is
effectively rank-1: one eigenvalue $\lambda_{\max}\approx 6$ with the
next-largest eigenvalue an order of magnitude smaller
($\lambda_{\max}/\lambda_2 \approx 20$), whereas in the uncoupled baseline
($g_{XY}=0$) the spectrum remains dispersed
($\lambda_{\max}/\lambda_2 \lesssim 1.3$). The observed ratio
($4.5$--$4.7$) is therefore not an artifact of the observable: it signals
genuine eigenvalue condensation driven by the coupling, separated from the
$g_{XY}=0$ baseline ($1.10\pm0.05$) by well over fifty standard deviations.

\item \textbf{Weak observer-dimension trend:}
Mean ratio decreases gently with observer dimension ($4.74 \to 4.49$ from
$d=2$ to $d=5$ at $n=30$), suggesting that additional observer degrees of
freedom dilute the directional coupling slightly. The trend is statistically
detectable (slope $\approx -0.08$ per unit $d$, $p=0.003$) but weak
($R^2=0.07$), and no single pair of dimensions differs significantly.
\end{enumerate}

%=============================================================================
\subsection{Connection to IKKT Matrix Model}

Our model shares the mechanism of directional symmetry breaking with the IKKT
matrix model~\cite{kim_nishimura_2011, kim_nishimura_2019}, in which Kim and
Nishimura found that spontaneous symmetry breaking
SO(9) $\rightarrow$ SO(3) requires sufficiently large $N$ for clear emergence
of (3+1)-dimensional spacetime. We note, however, that the corrected v6
results do \emph{not} exhibit the finite-size sensitivity that an earlier v5
analysis had suggested: within the tested range $N \leq 8$, our coupling
$g_{XY}=0.8$ produces robust emergence with no required $N$-dependent tuning.
Any comparison between the two models therefore rests on the shared
qualitative idea of competition between coupling energy and the phase space
of matrix configurations, not on a demonstrated finite-size crossover in our
model.

%=============================================================================
\section{Discussion}\label{sec:discussion}
%=============================================================================

\subsection{Advantages Over External Adapter}

The dynamical observer approach offers several advantages:

\begin{enumerate}[leftmargin=*]
\item \textbf{Physical Motivation:} Observer is dynamical system, not parameter.
\item \textbf{Emergent Direction:} Temporal direction emerges from observer dynamics.
\item \textbf{Robustness:} Works across $d = 2,3,4,5$ and all tested $N = 4$--$16$ with
      fixed coupling, no crossover and no required parameter tuning.
\item \textbf{Measurement Theory:} Observer as measurement device.
\end{enumerate}

\subsection{Relation to Philosophical Positions}

Our model is consistent with several philosophical positions within the 
restricted scope of our toy model:

\begin{itemize}[leftmargin=*]
\item \textbf{Leibniz (relational time):} Temporal direction created by observer.
\item \textbf{Barbour (timeless physics)~\cite{barbour_1999}:} Pre-observer state is symmetric.
\item \textbf{Rovelli (thermal time)~\cite{connes_rovelli_1994, rovelli_2018}:} Coupling creates preferred direction.
\item \textbf{Relational QM~\cite{rovelli_rqm_1996}:} States relative to observer.
\end{itemize}

\textbf{Important:} ``Consistent with'' does not mean ``proves'' or 
``validates.'' These philosophical positions remain open questions.

%=============================================================================
\section{Open Questions and Future Directions}\label{sec:open_questions}
%=============================================================================

Our results raise several important questions that warrant further investigation:

\begin{enumerate}[leftmargin=*]
\item \textbf{Gate saturation mechanism:}
At $N=8$ and $g_{XY} \gtrsim 1.05$ the aligned extent saturates against the
stability-gate ceiling, pinning the ratio near $10$. A fine scan of $g_{XY}$
in the interval $0.80$--$1.05$ (step $0.025$) shows the entry into saturation
is a \emph{smooth crossover}, not a sharp transition: the ratio rises
monotonically from $\approx 4.7$ to $\approx 10$, with the steepest single
step ($\Delta R \approx 1.7$) near $g_{XY}\approx 0.875$--$0.90$, and then
flattens onto the plateau for $g_{XY}\gtrsim 0.95$. Whether the saturation
value is a well-defined function of \texttt{max\_extent} remains open.

\item \textbf{Large-$N$ behavior:}
The corrected sampler shows robust emergence with fixed coupling up to
$N=16$ (Section~\ref{sec:fss}, Table~\ref{tab:fss_large_n}). Whether this
robustness persists at still larger $N$ ($N=18, 20, \ldots$) remains open;
a genuine finite-size crossover could still emerge beyond the currently
tested range, and would then need to be distinguished from the v5
artifact described in Section~\ref{sec:absence_crossover}.

\item \textbf{Observer-dimension dependence (dimensional dilution):}
With $n=30$ seeds per dimension, the mean ratio decreases monotonically from
$4.74$ ($d=2$) to $4.49$ ($d=5$). The trend is statistically detectable
(slope $\approx -0.08$ per unit $d$, $p=0.003$) but weak ($R^2=0.07$), and
the pairwise comparison between $d=3$ and $d=4$ is not significant
($p=0.61$). This is consistent with a mild ``dimensional dilution'' of the
directional coupling over additional observer directions, rather than any
sharp dependence. Confirming the physical origin of the dilution (and the
concomitant rise of CV with $d$, from $3.1\%$ at $d=2$ to $10.3\%$ at $d=5$)
would require still larger samples or an analytic argument.

\item \textbf{Universality:}
Would this mechanism work with different coupling structures or different
matrix models (e.g., BFSS, Lorentzian IKKT)? How universal is the
observer-induced temporalization phenomenon?

\item \textbf{Feedback coupling (back-reaction):}
The present model treats the observer $Y$ as an external field acting on $X$
with no back-reaction. To move beyond this one-way picture, one could
introduce a small back-reaction term from $X$ back to $Y$ (e.g.\ a
$g_{YX}$-weighted coupling in the observer action) and ask whether the
two-way coupled system still selects a stable temporal direction, or whether
the mutual feedback destabilizes the ordering. This is a natural step toward
a genuinely relational, rather than externally driven, notion of time.
\end{enumerate}

%=============================================================================
\section{Conclusion}\label{sec:conclusion}
%=============================================================================

We have presented a toy-level computational demonstration of temporal 
emergence from a dynamical observer. Our key findings:

\begin{enumerate}[leftmargin=*]
\item An observer system $Y$ can dynamically create the temporal direction 
      in a main system $X$ through directional coupling.
\item With the corrected v6 trajectory-mean sampler, we observe healthy
      temporal emergence (ratios $4.4$--$4.8$, alignment $100\%$) across the
      entire tested range: observer dimensions $d \in \{2,3,4,5\}$ and system
      sizes $N \in \{4,5,\ldots,16\}$, with fixed coupling $g_{XY}=0.8$.
\item No finite-size crossover or two-regime scaling is observed; the
      crossover and tuned-coupling picture reported in v5 was an artifact of
      a \texttt{step=2} sampling shortcut, now removed (Section~\ref{sec:absence_crossover}).
\item Perfect alignment ($X_{\max} = v_{\max}$) in all healthy emergence cases 
      confirms the observer's active role.
\item At strong coupling ($g_{XY} \gtrsim 1.05$, $N=8$) the aligned extent
      saturates against the stability gate, pinning the ratio near $10$; this
      is gate saturation, not a new phase.
\item Across observer dimensions, the mean ratio decreases weakly and
      monotonically with $d$ (slope $\approx -0.08$ per unit $d$, $p=0.003$,
      $R^2=0.07$ at $n=30$); the pairwise comparison between $d=3$ and $d=4$
      is not significant ($p=0.61$). The earlier $n=5$ ``marginally
      significant'' result ($p=0.049$) was a small-sample fluctuation and has
      been withdrawn.
\end{enumerate}

\textbf{Note on the v5 \textrightarrow{} v6 correction:} An earlier version of
this manuscript (v5) reported a finite-size crossover at $N\approx 7$, a
two-regime scaling picture with a tuned optimal coupling
$g_{XY}^{\mathrm{opt}}(N)$, and a non-monotonic variance spike at
$g_{XY}=1.05$. None of these are reproduced by the corrected v6 sampler. They
were traced to a \texttt{step=2} element-skip optimization in the v5 engine,
which for $N>6$ updated only even-indexed matrix elements and prevented full
thermalization. The v5 numbers have therefore been withdrawn and replaced by
the v6 trajectory-mean results reported here; the v5 manuscript is retained
unchanged as \texttt{manuscript\_v5.tex} for archival comparison.

\textbf{Appropriate Interpretation:} These results demonstrate that 
dynamical observer temporalization \textit{can be represented} in a 
consistent mathematical framework. They do not prove that time emerges 
from observer coupling in nature.

\textbf{Final Statement:} This work is a toy-level demonstration of 
mathematical consistency, not a discovery of physical laws. Following 
the philosophy of our exploratory notes, we emphasize that this is a 
\textbf{design framework demonstrating mathematical consistency}, 
not empirical physics.

\bibliographystyle{plain}
\bibliography{references}

\end{document}
